\documentclass[10pt,a4paper]{amsart}
\usepackage[margin=19mm]{geometry}
\usepackage{amsmath,amssymb,mathtools}
\usepackage[scr=rsfs,cal=euler]{mathalfa}
\usepackage{xfrac}
\usepackage[unicode,hidelinks,bookmarks=false]{hyperref}
\newcommand{\ie}{i.e.\ }
\newcommand{\cf}{cf.\ }
\newcommand{\resp}{respectively\ }
\newcommand{\Subsection}{\subsection}
\providecommand{\nobreakdash}{\nobreak\hbox{-}}
\setlength{\emergencystretch}{2em}
\multlinegap0pt
\usepackage[shortlabels]{enumitem}
\usepackage{amssymb}
\usepackage{mathtools}
\usepackage{xcolor}
\newtheorem{proposition}{Proposition}
\newtheorem{lemma}{Lemma}
\usepackage{cleveref}
\crefname{proposition}{Proposition}{Propositions}
\crefname{lemma}{Lemma}{Lemmas}
\newcommand{\F}{\mathbb F}
\newcommand{\Om}{\Omega}
\title{The small Davenport constant of the Heisenberg group of order $343$}
\author{Andreas Volkmann}
\date{Source: arXiv:2608.13747v1 (CC BY 4.0)}
\begin{document}
\maketitle

\noindent\emph{Source and licence.} The small Davenport constant of the Heisenberg group of order $343$. Version arXiv:2608.13747v1 (CC BY 4.0), distributed by arXiv under the Creative Commons Attribution 4.0 International licence, http://creativecommons.org/licenses/by/4.0/, as recorded in the arXiv licence metadata for this exact version and shown on the abstract page https://arxiv.org/abs/2608.13747v1. Full licence notice: \texttt{licenses/arxiv-2608.13747-CC-BY-4.0.txt}.\par\smallskip

\noindent\textit{Editorial scope.} Counted unit: the article's proposition prop:Ffacts (source0.tex lines 345--354). The article's complete original proof of it is delivered with it (source0.tex lines 356--387, one \texttt{proof} environment of 32 lines). No mathematics is written, repaired or re-derived: the statement and the proof are the article's own lines, in the article's own notation, and no retained line is substituted or re-flowed.  Retained uncounted as the source-local context the proof needs: lines 441--448.  Nothing was omitted from any retained span, so the package carries no omission record.  No retained line contains a citation, so the wrapper carries no bibliography and no bibliography entry.  No retained line contains a figure environment, an image-inclusion command or drawing code, so no image is delivered and the package declares no asset dependency.  Editorial formatting only, in addition: an AMS \texttt{amsart} preamble that restates the article's own declarations and macros used by the retained lines, namely the environments \texttt{proposition}, \texttt{lemma} on separate counters, together with the standard packages those lines need.  The retained environments print in this document as Proposition 1, Lemma 1 in order of appearance, whereas the article prints the same items with the numbering of its own full document.  The complete native source of the article as distributed in the arXiv e-print for this exact version is bundled unchanged as \texttt{sources/207609/source0.tex}.
\begin{lemma}[Exact cross-sum recursion]\label{lem:omega-rec}
The cross-sum set satisfies
\[
\Om(0)=\{0\},\qquad
\Om(n)=\bigcup_{v:n_v>0}
 \left(\Om(n-e_v)+(A(n)-a_v)b_v\right).
\]
\end{lemma}

\begin{proposition}[Finite two-line facts]\label{prop:Ffacts}
The following statements hold.
\begin{enumerate}[label=(F\arabic*)]
\item There are $47$ minimal zero-sum multisets over $C_7$. Their counts in lengths $2$ through $7$ are, respectively, $3$, $8$, $12$, $12$, $6$, and $6$. Among the $\binom{48}{2}=1128$ unordered pairs with repetition of minimal blocks on two independent directions, exactly $6$ pair--pair blocks are nonfull with $|\Om|=3$, exactly $18$ pair--triple blocks are nonfull with $|\Om|=5$, and all other blocks are full: $\Om=\F_7$.
\item A size-$7$ line multiset all of whose minimal zero-sums are pairs is supported on one of the three opposite pairs $\{1,6\},\{2,5\},\{3,4\}$.
\item For $P_i$ and $P_j$, with $i,j\geq1$, placed on the two independent lines, one has $|\Om|=3$ for $(i,j)=(1,1)$, $|\Om|=5$ for $(1,2)$ or $(2,1)$, and $\Om=\F_7$ otherwise.
\item Every triple--double-pair block, every $(\text{pair}\sqcup\text{triple})$--pair block, and every $(\text{triple}\sqcup\text{triple})$--pair block is full.
\item No size-$12$ line multiset contains a minimal zero-sum triple but no minimal zero-sum of length at least $4$. At size $11$, the only such multisets are the six scalar multiples of $\{1^6,5,6^4\}$; each contains a pair and a disjoint triple.
\end{enumerate}
\end{proposition}

\begin{proof}[Exhaustive verification]
A line multiset is a vector $c=(c_1,\ldots,c_6)\in\mathbb N^6$.
For each required size, the enumerator loops over all weak
compositions, tests $\sum_{t=1}^6tc_t\equiv0\pmod7$, and tests
minimality by looping over every vector $0\leq d\leq c$ other than
$0$ and $c$.  This produces the stated $47$ minimal blocks and length
distribution.  For every unordered pair with repetition, the block is
placed on the coordinate axes and its cross-sum mask is evaluated by
the exact recursion in \cref{lem:omega-rec}; its population count gives
(F1).

For (F3), the recursion is applied separately to the six basis blocks
$(P_i,P_j)$ with
\[
(i,j)\in\{(1,1),(1,2),(2,1),(2,2),(3,1),(1,3)\}.
\]
These give cross-sum cardinalities $3,5,5,7,7,7$, respectively.
Every remaining pair $(i,j)$ with $i,j\geq1$ contains one of the last
three full zero-sum blocks.  Ordering that subblock first and the
remaining terms in a fixed order shows that the larger block is full
as well.

For (F2) and (F5), the same weak-composition enumeration is run at the
stated total sizes.  The program records all contained minimal
zero-sum vectors and then applies the predicates in (F2) and (F5)
literally.  For (F4), it enumerates the indicated disjoint tuples of
minimal blocks, adds their count vectors, and again applies
\cref{lem:omega-rec}.  Every loop has an explicit finite domain and no
random or floating-point step.  A separately written enumerator
reproduced the counts.  The archived standalone verifier, its deterministic
output, and their hashes permit the displayed facts to be checked directly.
\end{proof}

\end{document}
